\section{Introduction}
Reinforcement learning (RL) is a central component of language-model post-training. It can optimize models against learned human-preference rewards~\citep{Ouyang2022} or verifiable task outcomes~\citep{shao2024deepseekmath,guo2025deepseek}. Proximal policy optimization (PPO; \citealp{schulman2017proximal}) provides a widely used training objective, but its actor--critic implementation requires a value model in addition to the policy. Training and storing this critic increases the resource requirements of large-model alignment.

Critic-free methods avoid this extra model by estimating advantages directly from sampled rewards. ReMax uses a greedy response as a baseline~\citep{li2023remax}, RLOO uses a leave-one-out baseline~\citep{2024Back}, and group relative policy optimization (GRPO) standardizes rewards within each prompt's response group~\citep{shao2024deepseekmath}. These methods differ in both their baselines and their normalization rules. In particular, RLOO does not inherently require the group-standard-deviation normalization used by GRPO.

This paper focuses on the \emph{scope of advantage normalization}. A prompt-level standard deviation is estimated from only a few responses and introduces a separate scale for every prompt. As a result, the update depends on the composition of each sampled group, as well as on its rewards. Homogeneous groups contribute no centered reward signal, whereas groups with different reward spreads are rescaled independently. These effects motivate separating two decisions: whether to use a prompt-specific baseline and whether to use a prompt-specific scale.

We introduce \ours, a critic-free framework that uses a shared normalization across the rollout batch. The basic variant combines Monte Carlo returns with PPO clipping and supports $k\geq 1$ responses per prompt. At $k=1$, it allocates each rollout to a different prompt. The \baseline variant retains a within-prompt mean baseline for $k>1$, then normalizes globally and applies a separate KL regularizer. This distinction preserves the option of comparing responses to the same prompt without assigning each prompt its own standard-deviation scale.

Our contributions are threefold. First, we specify both variants within a common clipped policy objective, including the loss reduction used in our experiments. Second, we characterize the finite-sample effects of self-normalization and establish convergence to population standardization under independent sampling. Third, we evaluate the methods across preference-based alignment, reasoning, and tool use. The results show competitive alignment scores with single-response sampling, improved held-out performance in several reasoning settings, and the highest average tool-use score among the compared methods.
